% sections/G_production.tex: Lane m4 (the cue-free production pipeline).
% Numbers only through \R / \Rci; figures through \genfig, tables through \gentab (report/README.md).
% One \subsection per decision group, so that each can be trimmed on its own.
\section{The cue-free production pipeline}\label{app:prod}

In production (the Stage~3 inference link) no GO cue marks the movement.
The person holds the phone at arm's length and brings it to the nose once or a few times.
The pipeline must find each approach in the continuous sensor stream, ignore the return and every other movement, decide when to stop, and output a calibrated \pright{}.
Lane m4 built this pipeline and simulated it on the recordings; this appendix states each of its decisions with the options, the evidence and the rationale.

\textbf{Which hand classifier.}
Every production number in this appendix uses lane m4's own reference hand classifier, yaw8: a mirror-augmented logistic regression on the cumulative yaw about gravity ($\vomega\cdot\ghat$) at 8 instants of the approach window.
It does not use the final model, the body-frame 1D CNN of Appendix~\ref{app:final}.
The hand classifier is a plug-in: the detector, the calibration and the stopping rule take any model that reads a window, and the final model inside this simulation is the first next step (Appendix~\ref{app:prod:caveats}).

\textbf{Unit, evaluation and scoring.}
The unit is one full recording of a primary trial: rest, approach, hold and return, without the app's stored event times.
There are \R{m4/n_recordings} recordings of \R{m4/n_persons} persons, evaluated leave-one-person-out; every fitted part (segmentation setting, segment classifier, hand model, calibrator, aggregation scale) is fitted on the training persons only.
An abstention (no approach accepted, or a calibrated \pright{} of exactly 0.5) is never a right-hand answer: it scores half a correct answer.
Next to this half-credit accuracy, each result gives the coverage (the share answered), the accuracy on the answered recordings and the forced-choice accuracy (the best candidate always scored).
In the tables of this appendix, `answered' means that an approach was accepted; it still counts the calibrated outputs of exactly 0.5 (Appendix~\ref{app:prod:calibration}), which are abstentions and score half.
For the production pipeline, the share of recordings with a hand shown is therefore \R{derived/m4_prod_shown_coverage} (\R{derived/m4_prod_n_shown} of \R{m4/n_recordings}), right on \R{derived/m4_prod_shown_acc} of those; these two numbers are exact arithmetic on the saved counts, done by the report's asset step.
Natural use is right-hand recall on the free-choice trials (all right-hand so far; a constant right answer would score 100~\%).

\textbf{Tests.}
Paired differences are means over persons with a person-bootstrap CI and an exact sign-flip test over persons.
Holm's adjustment runs within named families: \emph{trial} (the \R{m4/paired_family_n} comparisons of Table~\ref{tab:m4_paired}), \emph{stress} (\R{m4/stress_family_n} comparisons of the stress test), \emph{sessions} (\R{m4/sess_family_n} session comparisons), \emph{calibrators} (each calibrator against none) and \emph{grip} (\R{m4/grip_family_n} comparisons of the grip check).
The production pipeline was compared with \R{m4/n_configs_final} cue-free pipelines on the outer folds, so its own LOSO numbers are exploratory; the segmentation setting and the calibrator are chosen inside each fold.

\textbf{Headline.}
With one touch, the production pipeline reaches a half-credit accuracy of \Rci{m4/prod_acc}. It accepts an approach in \R{m4/prod_coverage} of the recordings (accuracy \R{m4/prod_acc_answered} on those, a 0.5 output scoring half) and shows a hand for \R{derived/m4_prod_shown_coverage}, right on \R{derived/m4_prod_shown_acc} of those; its forced-choice accuracy is \R{m4/prod_acc_forced}; its AUC is \Rci{m4/prod_auc}.
With the stopping rule on synthetic multi-touch sessions it reaches \Rci{m4/stop_1_acc} after \R{m4/stop_1_attempts} accepted attempts on average (exploratory; Appendix~\ref{app:prod:stopping}).

% -------------------------------------------------------------------------------------------------
\subsection{The baseline detector v0 and the evaluation contract}\label{app:prod:v0}

\textbf{Options.}
A fixed rule built from the app's own thresholds, or a detector learned from labelled segments.
The fixed rule, v0, comes first: it sets the baseline and the outcome categories that every later detector reports.

\textbf{The rule.}
v0 takes the first sustained movement (smoothed $|\vomega| > 30$~deg/s for 100~ms) that follows at least 1~s of stillness ($|\vomega| < 20$~deg/s), and ends it at the next 250~ms of stillness.
It reads the gyroscope only.
Adding an accelerometer condition ($|\va| < 0.5$~m/s$^2$) leaves \R{describe/still_both_share_1s} of the approaches with at least 1~s of stillness before them (median \R{describe/still_both_median_s}~s), against \R{describe/still_before_onset_share_1s} with the gyroscope alone (median \R{describe/still_before_onset_median_s}~s).

\textbf{Evidence.}
Figure~\ref{fig:app:prod-v0} and Table~\ref{tab:describe_cuefree_v0} score v0's first detection against the oracle window, the approach bounds the app stored, by their intersection over union (IoU).
The first detection is the approach (IoU $\ge 0.5$) in \R{describe/cuefree_share_approach} of the primary recordings, part of it in \R{describe/cuefree_share_approach_partial}, an earlier movement in \R{describe/cuefree_share_earlier_movement}, the return in \R{describe/cuefree_share_return} and nothing in \R{describe/cuefree_share_none}.
The approach is among all of v0's detections in \R{describe/cuefree_share_any_detection} of the recordings.
When v0 finds the approach, its onset is \R{describe/cuefree_onset_error_median_ms}~ms late (median; 90th percentile of the absolute error \R{describe/cuefree_onset_abs_error_p90_ms}~ms).
With yaw8 on its first detection, v0 reaches \Rci{m4/vzero_acc} (coverage \R{m4/vzero_coverage}, answered \R{m4/vzero_acc_answered}, forced choice \R{m4/vzero_acc_forced}).

\textbf{Decision and rationale.}
v0 is the baseline, and its categories (approach, part, earlier movement, return, nothing, rejected) are the contract of every detector in this appendix.
Its main failure is an earlier movement (the arm settling after the previous trial, a fidget), not the return.

\begin{figure}[htbp]
  \centering
  \genfig{core/S2_16_cuefree_v0}
  \caption{Baseline cue-free detector v0 on the primary recordings, by platform. Left: what v0 finds first (approach: IoU $\ge 0.5$ with the oracle window). Middle: onset error when it finds the approach. Right: stillness before the approach in v0's own signal, which the collection app enforced; the shares are therefore upper bounds for production (Figure~\ref{fig:app:prod-v0-stress}).}
  \label{fig:S2_16_cuefree_v0}\label{fig:app:prod-v0}
\end{figure}
\gentab{describe/describe_cuefree_v0}
\FloatBarrier

% -------------------------------------------------------------------------------------------------
\subsection{The stillness stress test: recorded shares are upper bounds}\label{app:prod:stress}

\textbf{Problem.}
The collection app enforced the stillness before every approach: it waited for 400~ms of stillness, then a 400--900~ms foreperiod, then the reaction time.
Production has no GO cue and no hold-still gate, so a user may start moving without holding still.
Every share measured on the recordings as they are is therefore an upper bound.

\textbf{Method.}
The stress test cuts the still run before each approach to a given length: the samples right after the preceding activity are removed and the rest of the recording moves back.
Every pipeline then runs again on the shortened recording, fitted without the held-out person.

\textbf{Evidence for v0.}
Figure~\ref{fig:app:prod-v0-stress} and Table~\ref{tab:describe_cuefree_stress} show that v0 needs its own 1~s rule.
It finds the approach first in \R{describe/stress_recorded_share_approach} of the recordings as recorded and in \R{describe/stress_keep1000ms_share_approach} with 1~s of stillness, but in \R{describe/stress_keep750ms_share_approach} with 0.75~s and \R{describe/stress_keep0ms_share_approach} with none.
With a random stillness, uniform between 0 and 1.5~s, it finds the approach first in \R{describe/stress_random_share_approach}.

\textbf{Evidence for the full pipelines.}
Figure~\ref{fig:app:prod-stress} and Table~\ref{tab:m4_stress} run the production pipeline, v0 and a detector whose stillness feature sees up to \R{m4/seg_still_cap_s}~s (context capped at 3~s) through the same cuts.
The production pipeline scores \R{m4/stress_production_rec_acc} as recorded, \R{m4/stress_production_1_acc} with 1~s, \R{m4/stress_production_0p9_acc} with 0.9~s, \R{m4/stress_production_0p75_acc} with 0.75~s, \R{m4/stress_production_0p5_acc} with 0.5~s and \R{m4/stress_production_0_acc} with none.
With no stillness it answers only \R{m4/stress_production_0_cov} of the recordings and scores the return in \R{m4/stress_production_0_ret}; v0 then answers \R{m4/stress_vzero_0_cov}, scores the return in \R{m4/stress_vzero_0_ret} and falls to \R{m4/stress_vzero_0_acc}.

\textbf{Options for the stillness feature.}
The stillness before a segment is a strong feature of the segment classifier (Appendix~\ref{app:prod:detector}).
It can see the whole stillness (capped at 3~s) or only up to the instructed 1~s.
Paired on the same recordings (robust detector minus the context capped at 3~s, both uncalibrated, family \emph{stress}), the difference is \Rb{m4/stresspair_robust_vs_full_rec_d}~[\Rb{m4/stresspair_robust_vs_full_rec_lo}, \Rb{m4/stresspair_robust_vs_full_rec_hi}]~points as recorded (Holm $p = \R{m4/stresspair_robust_vs_full_rec_p_holm}$), \Rb{m4/stresspair_robust_vs_full_1_d}~[\Rb{m4/stresspair_robust_vs_full_1_lo}, \Rb{m4/stresspair_robust_vs_full_1_hi}]~points with 1~s (Holm $p = \R{m4/stresspair_robust_vs_full_1_p_holm}$), \Rb{m4/stresspair_robust_vs_full_0p9_d}~[\Rb{m4/stresspair_robust_vs_full_0p9_lo}, \Rb{m4/stresspair_robust_vs_full_0p9_hi}]~points with 0.9~s (Holm $p = \R{m4/stresspair_robust_vs_full_0p9_p_holm}$) and \Rb{m4/stresspair_robust_vs_full_0_d}~[\Rb{m4/stresspair_robust_vs_full_0_lo}, \Rb{m4/stresspair_robust_vs_full_0_hi}]~points with none (Holm $p = \R{m4/stresspair_robust_vs_full_0_p_holm}$).
The production detector against v0 at 0.9~s: \Rb{m4/stresspair_production_vs_v0_0p9_d}~[\Rb{m4/stresspair_production_vs_v0_0p9_lo}, \Rb{m4/stresspair_production_vs_v0_0p9_hi}]~points (Holm $p < 0.001$).

\textbf{Decision and rationale.}
The stillness feature is capped at the instructed 1~s (the robust feature set), and the inference page must ask for at least 1~s of stillness and check it, measured exactly as the feature is (the causal activity below the stillness threshold).
The capped detector is better where production will operate, at the instructed 1~s; it is not better as recorded, where the long enforced stillness favours the uncapped one.
Caveat: the advantage appears only at and just below its own cap.

\begin{figure}[htbp]
  \centering
  \genfig{core/S2_19_cuefree_stress}
  \caption{Stress test of v0: the stillness before each approach cut to the given length, then v0 run again (right margin: as recorded). Below its own 1~s rule, v0 can no longer separate the approach from what came before.}
  \label{fig:S2_19_cuefree_stress}\label{fig:app:prod-v0-stress}
\end{figure}
\gentab{describe/describe_cuefree_stress}
\begin{figure}[htbp]
  \centering
  \genfig{m4/M4_03_stress}
  \caption{Stress test of the full pipelines. Left: half-credit accuracy (band: 95~\% CI). Right: share answered (solid) and share where the scored segment is the return (dotted). The production detector sees at most the instructed 1~s of stillness; with 1~s it scores \R{m4/stress_production_1_acc} against \R{m4/stress_uncapped_1_acc} for the context capped at 3~s. Supports capping the stillness feature at 1~s and asking for that stillness on the page.}
  \label{fig:M4_03_stress}\label{fig:app:prod-stress}
\end{figure}
{\renewcommand{\arraystretch}{0.9}\gentab{m4/m4_stress}}
\FloatBarrier

% -------------------------------------------------------------------------------------------------
\subsection{The return: a time-reversed approach}\label{app:prod:return}

\textbf{Problem.}
After the hold, the arm returns to arm's length.
The return retraces the approach backwards, so its chiral evidence is reversed.

\textbf{Evidence.}
Figure~\ref{fig:app:prod-return-yaw} and Table~\ref{tab:describe_return_risk} compare the class-mean yaw rate of the approach and of the return in the same \R{describe/return_trials} trials with a clean return window.
The sign rule fitted on approaches (right if and only if the first-half yaw is positive) is right on \R{describe/return_rule_acc_approach} of the approaches but on \R{describe/return_rule_acc_return} of the returns (\R{describe/return_rule_acc_return_net} on the net yaw of the return); the AUC of the first-half yaw on the return is \R{describe/return_auc_first_half}.
A detector that scores the return therefore answers the other hand.

\textbf{Decision.}
Returns are never pooled with approaches, and the detector must reject them: the segment classifier has a return class (Appendix~\ref{app:prod:detector}).

\textbf{Option: the return as evidence.}
The return carries the hand too, with the opposite sign.
Figure~\ref{fig:app:prod-return} tests a separate return head (the same hand model fitted on returns).
On the stored windows of the trials with a clean return, the approach model scores \Rci{m4/ret_app_acc}, the return head \Rci{m4/ret_head_acc} and the sum of their logits \Rci{m4/ret_sum_acc}; the sum minus the approach is \Rcib{m4/d_ret_sum_vs_app}~points (Holm $p = \R{m4/d_ret_sum_vs_app_p_holm}$, family \emph{trial}).
Cue-free, adding the return head when the segment after the chosen approach is classified as a return gives \R{m4/prod_return_acc} against \R{m4/robust_acc} (both uncalibrated, coverage unchanged; answered \R{m4/prod_return_acc_answered}, forced choice \R{m4/prod_return_acc_forced}); paired \Rcib{m4/d_return_vs_prod}~points (Holm $p = \R{m4/d_return_vs_prod_p_holm}$).
The price is time: the return's evidence arrives a median \R{m4/return_wait_ms}~ms after the end of the approach (90th percentile \R{m4/return_wait_p90_ms}~ms).

\textbf{Decision and rationale.}
The return head is not in production yet.
No accuracy difference is shown after adjustment, and the summed logits need their own cross-fitted calibration.
Between two touches the return happens anyway, so in a multi-touch session it would delay only the final answer; it is a next step.

\begin{figure}[htbp]
  \centering
  \genfig{core/S2_15_return_vs_approach}
  \caption{Class-mean yaw rate about gravity in the approach (left) and in the return (right) of the same \R{describe/return_trials} trials, by hand (band: 95~\% CI). The return mirrors the approach in time, so a sign rule fitted on approaches answers the other hand on returns.}
  \label{fig:S2_15_return_vs_approach}\label{fig:app:prod-return-yaw}
\end{figure}
\gentab{describe/describe_return_risk}
\begin{figure}[htbp]
  \centering
  \genfig{m4/M4_10_return}
  \caption{The return as evidence (half-credit accuracy, 95~\% CI). Top three: stored windows of the \R{m4/ret_app_n} trials with a clean return. Bottom two: the cue-free production detector (uncalibrated) without and with the return head, on all \R{m4/n_recordings} recordings.}
  \label{fig:M4_10_return}\label{fig:app:prod-return}
\end{figure}
\FloatBarrier

% -------------------------------------------------------------------------------------------------
\subsection{The learned detector: candidate segments and the segment classifier}\label{app:prod:detector}

\textbf{Candidate segments.}
A causal activity signal follows $|\vomega|$ through a one-pass second-order \R{m4/seg_smooth_hz}~Hz Butterworth low-pass; every bound is moved back by its group delay of \R{m4/causal_delay_ms}~ms.
Hysteresis cuts the stream into bouts: a bout becomes a movement once the activity exceeds \R{m4/seg_hi_dps}~deg/s for \R{m4/seg_sustain_ms}~ms, and it ends when the activity stays below a lower threshold for a sustain time.
Unlike v0, no stillness is required before a bout.
Three settings of the lower threshold and the sustain time are compared inside each fold; 15~deg/s with 250~ms was chosen in \R{m4/prod_seg_cfg_mode_folds} of \R{m4/prod_folds} folds.

\textbf{The segment classifier.}
A multinomial logistic regression (C = 1) labels every bout as an approach, a return or another movement, from hand-agnostic features scaled by their median and interquartile range and winsorised.
Its labels come from the app's events of the training persons.
Every feature is mirror-even (tested), so the segment classifier carries no hand information.
The candidate with the highest P(approach) is scored; it is rejected (an abstention) when P(approach) is below \R{m4/det_threshold}.
A segment is confirmed \R{m4/lookahead_s}~s after its end, when its window is complete.

\textbf{Evidence: what separates the approach from the return.}
Figure~\ref{fig:app:prod-segfeat} shows the four strongest features.
The displacement along the screen normal separates approaches from returns with an AUC of \R{m4/auc_disp_z0_vs_return} (\R{m4/auc_disp_z0_vs_other} against other movements).
The change of the up vector's screen-normal component has an AUC of \R{m4/auc_d_up_z_vs_return}, the reverse ordering.
The vertical displacement reaches \R{m4/auc_disp_vert_vs_return}, and the stillness before the segment, capped at 1~s, reaches \R{m4/auc_still_before_inst_vs_other} against other movements (\R{m4/auc_still_before_inst_vs_return} against returns).
The screen-normal features are in the device frame, so they change sign when the screen faces away (Appendix~\ref{app:prod:grip}).

\textbf{Options: five feature sets.}
\emph{physics} (the movement features); \emph{full} (physics plus the stillness before, capped at 3~s, and recording context); \emph{nested} (physics or full chosen inside each fold: full in \R{m4/nested_features_chosen_mode_folds} folds); \emph{robust} (physics plus the stillness before, capped at the instructed 1~s, and the stillness after); and \emph{grip} (robust with the four screen-normal quantities signed by the side of the phone that faces up over the segment).
Table~\ref{tab:m4_models} gives physics \R{m4/physics_acc}, full \R{m4/full_acc}, robust \R{m4/robust_acc} and grip \R{m4/grip_acc}.
Paired (family \emph{trial}): robust minus nested \Rcib{m4/d_robust_vs_nested}~points (Holm $p = \R{m4/d_robust_vs_nested_p_holm}$); grip minus robust \Rcib{m4/d_grip_vs_robust}~points (Holm $p = \R{m4/d_grip_vs_robust_p_holm}$).
The grip set is invariant to the screen-away grip by construction, but the side that faces up is a poor grip estimate on real recordings: its sign is negative for \R{m4/grip_sign_neg_approach} of the approach and \R{m4/grip_sign_neg_return} of the return segments of screen-toward recordings.

\textbf{Decision and rationale.}
The robust set is the production detector, an exploratory choice among \R{m4/n_configs_detector} feature sets.
It is the set that does not reward the enforced stillness of the protocol (Appendix~\ref{app:prod:stress}); the screen-away grip is handled by the instructions and the page's grip check instead of the grip set.

\textbf{What the detector scores.}
Table~\ref{tab:m4_detection} scores each detector's chosen segment against the stored events.
The production detector scores the approach in \R{m4/prod_share_approach} of the recordings, an earlier movement in \R{m4/prod_share_earlier_movement} and the return in \R{m4/prod_share_return}; it rejects \R{m4/prod_share_rejected} instead of scoring a wrong segment.
v0 scores the approach in \R{describe/cuefree_share_approach} and an earlier movement in \R{describe/cuefree_share_earlier_movement}.
When the approach is found, the median bound errors (detected minus stored) are \R{m4/prod_onset_err_ms}~ms at the onset and \R{m4/prod_end_err_ms}~ms at the end.

\begin{figure}[htbp]
  \centering
  \genfig{m4/M4_11_segment_features}
  \caption{Hand-agnostic features of the candidate segments of the recordings, by their label from the app's events (densities; AUC against return and against other movements in the panel titles). They separate the approach from the return without any hand information and are the basis of the segment classifier.}
  \label{fig:M4_11_segment_features}\label{fig:app:prod-segfeat}
\end{figure}
\gentab{m4/m4_detection}
\FloatBarrier

% -------------------------------------------------------------------------------------------------
\subsection{Pipelines, references and data-treatment ablations}\label{app:prod:pipelines}

\textbf{The ladder.}
Table~\ref{tab:m4_models} lists every cue-free pipeline and its references: the same hand model on the oracle window (no detection), v0, the five feature sets, the calibrators, the return head and the ablations.
Table~\ref{tab:app:prod-scoring} adds, for the main rungs, what an abstention hides: coverage, accuracy on the answered recordings and forced-choice accuracy, on the primary recordings and on natural use.

\begin{table}[htbp]
  \centering\small
  \caption{Scoring of the main cue-free pipelines (yaw8 hand model; LOSO, \R{m4/n_recordings} recordings of \R{m4/n_persons} persons). Accuracy: an abstention scores half credit. Answered: coverage. Forced: the best candidate always scored. Natural use: right-hand recall on the \R{m4/prod_nat_n} free-choice trials (all right-hand so far; a constant right answer would score 100~\%), as answered of \R{m4/prod_nat_n}, recall on the answered, forced choice and half credit.}
  \label{tab:app:prod-scoring}
  \fitwidth{%
  \begin{tabular}{lrlrrrrrrr}
    \toprule
    & \multicolumn{5}{c}{Primary recordings} & \multicolumn{4}{c}{Natural use (right-hand recall)} \\
    \cmidrule(lr){2-6}\cmidrule(lr){7-10}
    Pipeline & Accuracy & 95~\% CI & Answered & Acc.\ answered & Forced & Answered & Answered & Forced & Half credit \\
    \midrule
    oracle window & \R{m4/oracle_acc} & [\Rb{m4/oracle_acc_lo}, \Rb{m4/oracle_acc_hi}] & \R{m4/oracle_coverage} & \R{m4/oracle_acc_answered} & \R{m4/oracle_acc_forced} & \R{m4/oracle_nat_answered} & \R{m4/oracle_nat_acc_answered} & \R{m4/oracle_nat_acc_forced} & \R{m4/oracle_nat_acc} \\
    v0 (first movement) & \R{m4/vzero_acc} & [\Rb{m4/vzero_acc_lo}, \Rb{m4/vzero_acc_hi}] & \R{m4/vzero_coverage} & \R{m4/vzero_acc_answered} & \R{m4/vzero_acc_forced} & \R{m4/vzero_nat_answered} & \R{m4/vzero_nat_acc_answered} & \R{m4/vzero_nat_acc_forced} & \R{m4/vzero_nat_acc} \\
    learned, physics & \R{m4/physics_acc} & [\Rb{m4/physics_acc_lo}, \Rb{m4/physics_acc_hi}] & \R{m4/physics_coverage} & \R{m4/physics_acc_answered} & \R{m4/physics_acc_forced} & \R{m4/physics_nat_answered} & \R{m4/physics_nat_acc_answered} & \R{m4/physics_nat_acc_forced} & \R{m4/physics_nat_acc} \\
    learned, full & \R{m4/full_acc} & [\Rb{m4/full_acc_lo}, \Rb{m4/full_acc_hi}] & \R{m4/full_coverage} & \R{m4/full_acc_answered} & \R{m4/full_acc_forced} & \R{m4/full_nat_answered} & \R{m4/full_nat_acc_answered} & \R{m4/full_nat_acc_forced} & \R{m4/full_nat_acc} \\
    learned, robust (uncal.) & \R{m4/robust_acc} & [\Rb{m4/robust_acc_lo}, \Rb{m4/robust_acc_hi}] & \R{m4/robust_coverage} & \R{m4/robust_acc_answered} & \R{m4/robust_acc_forced} & \R{m4/robust_nat_answered} & \R{m4/robust_nat_acc_answered} & \R{m4/robust_nat_acc_forced} & \R{m4/robust_nat_acc} \\
    learned, grip & \R{m4/grip_acc} & [\Rb{m4/grip_acc_lo}, \Rb{m4/grip_acc_hi}] & \R{m4/grip_coverage} & \R{m4/grip_acc_answered} & \R{m4/grip_acc_forced} & \R{m4/grip_nat_answered} & \R{m4/grip_nat_acc_answered} & \R{m4/grip_nat_acc_forced} & \R{m4/grip_nat_acc} \\
    production (calibrated) & \R{m4/prod_acc} & [\Rb{m4/prod_acc_lo}, \Rb{m4/prod_acc_hi}] & \R{m4/prod_coverage} & \R{m4/prod_acc_answered} & \R{m4/prod_acc_forced} & \R{m4/prod_nat_answered} & \R{m4/prod_nat_acc_answered} & \R{m4/prod_nat_acc_forced} & \R{m4/prod_nat_acc} \\
    \bottomrule
  \end{tabular}}
\end{table}

\textbf{Detection loss.}
Production minus the oracle window is \Rcib{m4/d_prod_vs_oracle}~points (Holm $p = \R{m4/d_prod_vs_oracle_p_holm}$, family \emph{trial}; Table~\ref{tab:m4_paired}): this is the price of finding the movement.
Production minus v0 is \Rcib{m4/d_prod_vs_v0}~points (Holm $p = \R{m4/d_prod_vs_v0_p_holm}$: no accuracy difference shown) and \Rcib{m4/d_prod_vs_v0_ll} of log loss (Holm $p = \R{m4/d_prod_vs_v0_ll_p_holm}$).
The learned detector is kept because its probabilities are better and because v0 fails when the stillness is short (Appendix~\ref{app:prod:stress}).

\textbf{Training windows of the hand model.}
Options: the oracle windows or the detector's windows.
Detector minus oracle windows: \Rcib{m4/d_detwin_vs_oraclewin}~points (Holm $p = \R{m4/d_detwin_vs_oraclewin_p_holm}$).
Decision: oracle windows, the cached training arrays, so that any window model plugs in unchanged.

\textbf{Activity filter and stream low-pass.}
The activity signal runs causal (zero-phase minus causal: \Rcib{m4/d_zp_activity_vs_prod}~points, Holm $p = \R{m4/d_zp_activity_vs_prod_p_holm}$), because it must run online.
The stream itself is low-passed zero-phase over the buffer received so far, the code path of the training windows; a fully causal low-pass minus this is \Rcib{m4/d_causal_lowpass_vs_prod}~points (Holm $p = \R{m4/d_causal_lowpass_vs_prod_p_holm}$).
A replay re-standardised each raw record up to each decision time, as the page would: \R{m4/replay_matched} of the segments are found again, decisions change in \R{m4/replay_changed} of the recordings (statuses in \R{m4/replay_status_changed}), and the largest change of \pright{} is \R{m4/replay_max_dp}.

\textbf{Hand model.}
The second mirror-augmented hand model, chiral16, minus yaw8 is \Rcib{m4/d_chiral16_vs_yaw8}~points cue-free (Holm $p = \R{m4/d_chiral16_vs_yaw8_p_holm}$) and \Rcib{m4/d_oracle_chiral16_vs_yaw8}~points on the oracle window (Holm $p = \R{m4/d_oracle_chiral16_vs_yaw8_p_holm}$).
Decision: yaw8 stays the pre-specified default; the hand model is pluggable.

\textbf{Ablations.}
Table~\ref{tab:m4_ablations} varies one data treatment at a time against the uncalibrated production detector: feature scaling (z-score \Rcib{m4/d_abl_scale_z_vs_base}, none \Rcib{m4/d_abl_scale_none_vs_base}~points against median/IQR), the segment classifier's C nested in the folds (\Rcib{m4/d_abl_segc_nested_vs_base}~points) and the look-ahead (0.3~s \Rcib{m4/d_abl_look0p3_vs_base}, 0.8~s \Rcib{m4/d_abl_look0p8_vs_base}~points against 0.5~s).
None shows an accuracy difference after Holm's adjustment (family \emph{trial}).
The pre-specified values stay: median/IQR scaling, C = 1, and a 0.5~s look-ahead, which adds directly to the latency.

\textbf{P(approach) threshold.}
Re-thresholding the saved predictions on the outer folds (exploratory, outside the Holm family) trades coverage for accuracy: 0.30 gives \R{m4/thr_0p30_acc} (answers \R{m4/thr_0p30_coverage}; paired \Rcib{m4/thr_0p30_d}~points), 0.40 \R{m4/thr_0p40_acc} (\R{m4/thr_0p40_coverage}), 0.50 \R{m4/thr_0p50_acc} (\R{m4/thr_0p50_coverage}), 0.60 \R{m4/thr_0p60_acc} (\R{m4/thr_0p60_coverage}) and 0.70 \R{m4/thr_0p70_acc} (\R{m4/thr_0p70_coverage}).
The pre-specified 0.5 stays; a threshold chosen inside the folds is a candidate for the next version.

\gentab{m4/m4_models}
\gentab{m4/m4_paired}
\gentab{m4/m4_ablations}
\FloatBarrier

% -------------------------------------------------------------------------------------------------
\subsection{Calibration of one touch}\label{app:prod:calibration}

\textbf{Problem and detection.}
The product is the percentage shown to the user, so it must be calibrated.
Uncalibrated, the pipeline's expected calibration error is \R{m4/cal_none_ece}, above the range a perfectly calibrated \pright{} with the same values would show (null band \R{m4/cal_none_ece_null_lo}--\R{m4/cal_none_ece_null_hi}, $p = \R{m4/cal_none_ece_p_null}$); its calibration slope is \Rci{m4/cal_none_slope} and its intercept \Rci{m4/cal_none_intercept} (Figure~\ref{fig:app:prod-cal}, left).

\textbf{Options.}
Every calibrator is fitted on cross-fitted predictions of the training persons (person-grouped inner folds) and applied to the held-out person: none, temperature scaling \cite{guo2017calibration}, Platt scaling \cite{platt1999probabilistic}, isotonic regression, beta calibration and a symmetrised isotonic regression.
The hand model is mirror-augmented, so a mirror-image recording gets exactly $1 - \pright$; a calibrator keeps this only if it is antisymmetric, $P(-z) = 1 - P(z)$ for a logit $z$.
Platt, plain isotonic and beta are not: their largest deviation $|P(z) + P(-z) - 1|$ is \R{m4/cal_platt_mirror_dev_max}, \R{m4/cal_isotonic_mirror_dev_max} and \R{m4/cal_beta_mirror_dev_max}.

\textbf{Evidence.}
Table~\ref{tab:m4_calibration} and Figure~\ref{fig:app:prod-cal} (right) give each calibrator's paired log-loss difference to none (family \emph{calibrators}, Holm).
The log-loss difference to none is \Rcib{m4/cal_isotonic_d_ll} for plain isotonic (Holm $p = \R{m4/cal_isotonic_d_ll_p_holm}$) and \Rcib{m4/cal_isotonic_sym_d_ll} for the symmetrised isotonic (Holm $p = \R{m4/cal_isotonic_sym_d_ll_p_holm}$); temperature and Platt change little.
No calibrator shows a log-loss difference after adjustment; the choice rests on the calibration diagnostics below.

\textbf{Decision and rationale.}
The calibrator is chosen inside each fold among the antisymmetric ones (none, temperature, symmetrised isotonic) by person-grouped cross-validation of the cross-fitted logits; the symmetrised isotonic won in \R{m4/prod_calibrator_isotonic_sym_folds} of \R{m4/prod_folds} folds.
The production pipeline's ECE is \R{m4/cal_nested_ece} (null band \R{m4/cal_nested_ece_null_lo}--\R{m4/cal_nested_ece_null_hi}), its slope \Rci{m4/cal_nested_slope} and its intercept \Rci{m4/cal_nested_intercept}: consistent with calibration.
Against the uncalibrated pipeline its log loss changes by \Rcib{m4/cal_nested_d_ll} (Holm $p = \R{m4/cal_nested_d_ll_p_holm}$), its Brier score \cite{brier1950verification} by \Rcib{m4/cal_nested_d_brier} and its accuracy by \Rcib{m4/cal_nested_d_acc}~points.
The flat centre of the symmetrised isotonic map returns exactly 0.5 for \R{m4/prod_n_answered_tie} of the \R{m4/prod_n_answered} answered recordings: these show no hand and score half credit.
Per platform, Android has ECE \R{m4/cal_nested_android_ece} (null band \R{m4/cal_nested_android_ece_null_lo}--\R{m4/cal_nested_android_ece_null_hi}) and slope \Rci{m4/cal_nested_android_slope}; iOS has ECE \R{m4/cal_nested_ios_ece} (null band \R{m4/cal_nested_ios_ece_null_lo}--\R{m4/cal_nested_ios_ece_null_hi}) and slope \Rci{m4/cal_nested_ios_slope} (per-platform intercepts were not saved).

\begin{figure}[htbp]
  \centering
  \genfig{m4/M4_05_calibration}
  \caption{Calibration of one touch (10 equal-mass bins; band: 95~\% CI of the observed share; text: ECE with its perfect-calibration null band, slope with its CI). Left: uncalibrated, not consistent with calibration. Middle: production, the calibrator chosen inside each fold (symmetrised isotonic in \R{m4/prod_calibrator_isotonic_sym_folds} of \R{m4/prod_folds} folds), consistent. Right: paired log-loss difference of each calibrator to none (blue: antisymmetric).}
  \label{fig:M4_05_calibration}\label{fig:app:prod-cal}
\end{figure}
\gentab{m4/m4_calibration}
\FloatBarrier

% -------------------------------------------------------------------------------------------------
\subsection{The single-touch abstain rule and the decision latency}\label{app:prod:abstain}

\textbf{Abstain rule.}
A single-touch page can ask the user to repeat instead of answering.
Two rules were compared (Figure~\ref{fig:app:prod-coverage}): answer only when $\max(\pright, 1-\pright) \ge c$, or only when the best candidate's P(approach) reaches a threshold.
The confidence rule buys accuracy; the P(approach) rule barely does.
With $c = \R{m4/abstain_tau}$ the pipeline answers \R{m4/abstain_coverage} of the recordings and is right on \R{m4/abstain_acc} of those; its forced-choice accuracy stays \R{m4/prod_acc_forced}.
The same rule on the oracle windows lies above the cue-free curve at every coverage: part of the gap is detection.
The threshold is a product choice read off this curve, not tuned inside the folds; multi-touch sessions use the stopping rule instead (Appendix~\ref{app:prod:stopping}).

\textbf{Decision latency.}
The end of an approach is decided at the first stillness of the sustain time at the face (the deceleration to rest; the contact transient is not needed), plus the \R{m4/lookahead_s}~s look-ahead.
Figure~\ref{fig:app:prod-latency} shows the result: the decision comes a median \R{m4/lat_end_ms}~ms after the stored end of the approach (90th percentile \R{m4/lat_end_p90_ms}~ms) and \R{m4/lat_touch_ms}~ms after the first touch of the screen on the nose (90th percentile \R{m4/lat_touch_p90_ms}~ms).
A shorter look-ahead would cut the latency without an accuracy difference shown (Appendix~\ref{app:prod:pipelines}).

\begin{figure}[htbp]
  \centering
  \begin{subfigure}[t]{0.49\linewidth}
    \centering
    \genfig{m4/M4_07_coverage}
    \caption{Accuracy on the answered recordings against coverage, for a confidence threshold (blue), a P(approach) threshold (orange) and the confidence threshold on the oracle windows (grey).}
    \label{fig:M4_07_coverage}\label{fig:app:prod-coverage}
  \end{subfigure}\hfill
  \begin{subfigure}[t]{0.49\linewidth}
    \centering
    \genfig{m4/M4_09_latency}
    \caption{Decision time of the online pipeline after the stored end of the approach (blue) and after the first touch (yellow); recordings whose approach was found; values beyond the axis drawn at its edge.}
    \label{fig:M4_09_latency}\label{fig:app:prod-latency}
  \end{subfigure}
  \caption{The single-touch abstain rule (a) and the latency of the end-of-approach decision (b), production pipeline.}
  \label{fig:app:prod-coverage-latency}
\end{figure}
\FloatBarrier

% -------------------------------------------------------------------------------------------------
\subsection{Several touches: synthetic sessions and the k curve}\label{app:prod:sessions}

\textbf{Method.}
The recordings contain one approach each, so multi-touch sessions are simulated.
Each synthetic session concatenates \R{m4/sess_recordings} recordings of one person and hand, with the stillness at every junction cut to a given length; there are \R{m4/sess_n} sessions per setting.
Each fold's fitted pipeline runs over the held-out person's sessions online.
Every bout that the segment classifier accepts as an approach is an attempt, and the attempts' calibrated logits $z_i$ are combined:
\begin{equation}
  P_k = \sigma\!\left(s_k\,\bar z_k\right) = \sigma\!\left(\frac{1}{\tauk}\sum_{i=1}^{k} z_i\right), \qquad \tauk = \frac{k}{s_k},
  \label{eq:prod-aggregate}
\end{equation}
where $\sigma$ is the logistic function and $\bar z_k$ the mean of the first $k$ logits.

\textbf{Evidence: the k curve.}
Figure~\ref{fig:app:prod-attempts} (left) and Table~\ref{tab:m4_sessions_k} give the half-credit accuracy of the sign of the mean logit of up to $k$ accepted attempts.
With 1~s at each junction it is \Rci{m4/sess_1_k1_acc} for $k = 1$, \R{m4/sess_1_k2_acc} for $k = 2$, \Rci{m4/sess_1_k3_acc} for $k = 3$ and \Rci{m4/sess_1_k5_acc} for $k = 5$; \R{m4/sess_1_k5_avail} of the sessions reach 5 accepted attempts.
Keeping only the attempts that are true approaches (the attainable curve) gives \R{m4/sesstrue_1_k1_acc}, \Rci{m4/sesstrue_1_k3_acc} and \R{m4/sesstrue_1_k5_acc} for $k = 1$, 3 and 5.
With no stillness at the junctions the curve stays near chance (\R{m4/sess_0_k3_acc} at $k = 3$); with 0.5~s it reaches \R{m4/sess_0p5_k3_acc} and with 2~s \R{m4/sess_2_k3_acc}.

\textbf{Cause: false attempts.}
Only \R{m4/sessdet_precision} of the \R{m4/sessdet_attempts} accepted attempts are (part of) an approach.
The others are the arm settling before the approach (\R{m4/kind_other_before_share} of the attempts, \R{m4/sessdet_false_other_before}), returns (\R{m4/kind_return_share}, \R{m4/sessdet_false_return}) and later movements (\R{m4/kind_other_after_share}, \R{m4/sessdet_false_other_after}).
One attempt is right \R{m4/kind_approach_acc} of the time when it is the approach, \R{m4/kind_approach_partial_acc} when it is part of it, \R{m4/kind_other_before_acc} when it is the arm settling, \R{m4/kind_return_acc} when it is a return and \R{m4/kind_other_after_acc} for later movements.
The false attempts cancel the gain of more touches; rejecting them is the first next step.

\textbf{Option: the attempt rule.}
The pairing rule treats the bout after an accepted approach as its return and skips it; the plain rule (no pairing) takes every accepted bout.
Paired on the same sessions with the stopping rule (family \emph{sessions}), no pairing minus pairing is \Rcib{m4/sp_rule_no_pairing_vs_pairing_scale_sessions_d_acc}~points (Holm $p = \R{m4/sp_rule_no_pairing_vs_pairing_scale_sessions_p_holm}$, \R{m4/sp_rule_no_pairing_vs_pairing_scale_sessions_discordant} discordant sessions).
No pairing recalls \R{m4/sessdet_recall} of the true approaches against \R{m4/sessdetalt_recall} with pairing, at the same precision (\R{m4/sessdet_precision} and \R{m4/sessdetalt_precision}).
Decision: no pairing, because the pairing rule discards a true approach whenever an arm-settling movement before it was accepted.
For $k = 3$ the pairing rule scores \Rci{m4/sesspair_1_k3_acc}; the attempt rule is one of the exploratory session choices (Appendix~\ref{app:prod:stopping}).

\textbf{Not a cue-free number.}
The harness's $k = 3$ of the same pipeline, \Rci{m4/prod_acc_k3_indep}, averages three independent recordings of one person and hand, each with its best candidate.
It is not a cue-free session and is not comparable with the oracle-window $k = 3$ of other lanes.

\begin{figure}[htbp]
  \centering
  \genfig{m4/M4_04_attempts}
  \caption{Synthetic multi-attempt sessions (\R{m4/sess_n} per stillness setting). Left: half-credit accuracy of up to $k$ accepted attempts, by junction stillness (band: 95~\% CI at 1~s; labels: sessions with at least $k$ attempts); dark dashed: true approaches only; grey dotted: the pairing rule. Right: the stopping rule at 1~s, one point per $p_{\mathrm{stop}}$, for three aggregations. The gap between the 1~s curve and the true-approach curve is the cost of false attempts.}
  \label{fig:M4_04_attempts}\label{fig:app:prod-attempts}
\end{figure}
\gentab{m4/m4_sessions_k}
\FloatBarrier

% -------------------------------------------------------------------------------------------------
\subsection{The stopping rule and the calibration of several touches}\label{app:prod:stopping}

\textbf{The rule.}
After each accepted attempt, $P_k$ of Equation~\eqref{eq:prod-aggregate} is updated; the pipeline stops when $P_k$ leaves $[1 - p_{\mathrm{stop}}, p_{\mathrm{stop}}]$ or after \R{m4/k_max} attempts.
$p_{\mathrm{stop}} = \R{m4/p_stop}$ is a product choice fixed before the sessions were evaluated.

\textbf{Problem: the scale of several touches.}
Accepted attempts are not independent evidence: some are not approaches, and the attempts of one session share the person's bias.
Three aggregations were compared (Figure~\ref{fig:app:prod-k3cal}, Table~\ref{tab:m4_stop_calibration}): the mean logit ($s_k = 1$, nothing fitted); $s_k$ fitted on random recordings of the training persons (median $s_3 = \R{m4/scale_recordings_s3}$), which assumes neither problem; and $s_k = s_1 k^b$ fitted on cross-fitted synthetic sessions of the training persons (median $b = \R{m4/scale_sessions_b}$, $s_3 = \R{m4/scale_sessions_s3}$).
The sum of the logits (as if independent) is shown for $k = 3$ only.

\textbf{Evidence at the stop} (1~s junctions, $p_{\mathrm{stop}} = 0.90$).
Fitted on recordings: accuracy \R{m4/stopcal_scale_recordings_accuracy_half} with \R{m4/stopcal_scale_recordings_attempts} attempts, ECE \R{m4/stopcal_scale_recordings_ece} (null band \R{m4/stopcal_scale_recordings_ece_null_lo}--\R{m4/stopcal_scale_recordings_ece_null_hi}), slope \Rci{m4/stopcal_scale_recordings_slope}, intercept \R{m4/stopcal_scale_recordings_intercept}: not consistent with calibration.
Fitted on sessions: accuracy \R{m4/stopcal_scale_sessions_accuracy_half} with \R{m4/stopcal_scale_sessions_attempts} attempts, ECE \R{m4/stopcal_scale_sessions_ece} (null band \R{m4/stopcal_scale_sessions_ece_null_lo}--\R{m4/stopcal_scale_sessions_ece_null_hi}, at its upper edge), slope \Rci{m4/stopcal_scale_sessions_slope}, intercept \R{m4/stopcal_scale_sessions_intercept}: consistent with calibration.
Mean logit: accuracy \R{m4/stopcal_mean_accuracy_half} with \R{m4/stopcal_mean_attempts} attempts, ECE \R{m4/stopcal_mean_ece} (null band \R{m4/stopcal_mean_ece_null_lo}--\R{m4/stopcal_mean_ece_null_hi}), slope \Rci{m4/stopcal_mean_slope}, intercept \R{m4/stopcal_mean_intercept}: not consistent.
The intercepts' CIs were not saved.
Paired on the same sessions (family \emph{sessions}), recording-fitted minus session-fitted is \Rcib{m4/sp_agg_scale_recordings_vs_scale_sessions_d_acc}~points of accuracy and \Rcib{m4/sp_agg_scale_recordings_vs_scale_sessions_d_ll} of log loss (Holm $p = \R{m4/sp_agg_scale_recordings_vs_scale_sessions_ll_p_holm}$); session-fitted minus mean logit is \Rcib{m4/sp_agg_scale_sessions_vs_mean_d_acc}~points and \Rcib{m4/sp_agg_scale_sessions_vs_mean_d_ll} of log loss (Holm $p = \R{m4/sp_agg_scale_sessions_vs_mean_ll_p_holm}$).
No paired difference is shown.

\textbf{Decision and rationale.}
A declared rule picks the aggregation: among those consistent with calibration, the one with the fewest fitted parameters.
It picks $s_k$ fitted on sessions.
This choice is post hoc and exploratory: the session-fitted scale was added after the held-out sessions had shown the problem, and \R{m4/n_configs_sessions} configurations were compared (attempt rule $\times$ aggregation).

\textbf{Result.}
With 1~s junctions the stopping rule reaches \Rci{m4/stop_1_acc} with \R{m4/stop_1_attempts} accepted attempts and \R{m4/stop_1_touches} true touches by the decision on average; AUC \R{m4/stopcal_scale_sessions_roc_auc}, log loss \R{m4/stopcal_scale_sessions_log_loss}.
Every session answers (\R{m4/stop_1_share_no_attempt} without an attempt); \R{m4/stop_1_conf_share} stop on the threshold and are right \Rci{m4/stopcal_scale_sessions_conf_acc} of the time while claiming \R{m4/stopcal_scale_sessions_conf_claimed} on average; the other \R{m4/stop_1_share_k_max} stop at \R{m4/k_max} attempts.
With 0.5~s junctions it falls to \Rci{m4/stop_0p5_acc} and with none to \Rci{m4/stop_0_acc}.
Per platform: Android \R{m4/stopcal_android_acc} (ECE \R{m4/stopcal_android_ece}, null band \R{m4/stopcal_android_ece_null_lo}--\R{m4/stopcal_android_ece_null_hi}, slope \R{m4/stopcal_android_slope}; \R{m4/stopcal_android_n} sessions of \R{m4/stopcal_android_persons} persons), iOS \R{m4/stopcal_ios_acc} (ECE \R{m4/stopcal_ios_ece}, null band \R{m4/stopcal_ios_ece_null_lo}--\R{m4/stopcal_ios_ece_null_hi}, slope \R{m4/stopcal_ios_slope}; \R{m4/stopcal_ios_n} sessions of \R{m4/stopcal_ios_persons} persons); the CIs of these slopes and the intercepts were not saved, so these are descriptive.

\textbf{Exploratory readings of the same sessions.}
Table~\ref{tab:m4_stopping} sets touches against certainty: $p_{\mathrm{stop}} = 0.75$ gives \R{m4/pstop_0p75_acc} with \R{m4/pstop_0p75_attempts} attempts, 0.85 gives \R{m4/pstop_0p85_acc} with \R{m4/pstop_0p85_attempts}, 0.95 gives \R{m4/pstop_0p95_acc} with \R{m4/pstop_0p95_attempts}.
An abstain rule after \R{m4/k_max} attempts (answer only when $\max(P, 1-P) \ge p_{\mathrm{answer}}$) answers \R{m4/answer_0p60_coverage} of the sessions with \R{m4/answer_0p60_acc} right at $p_{\mathrm{answer}} = 0.6$, \R{m4/answer_0p70_coverage} with \R{m4/answer_0p70_acc} at 0.7 and \R{m4/answer_0p80_coverage} with \R{m4/answer_0p80_acc} at 0.8 (accuracy on the answered sessions).

\begin{figure}[htbp]
  \centering
  \genfig{m4/M4_06_k3_calibration}
  \caption{Calibration of several touches on the synthetic sessions (1~s junctions, production attempt rule; equal-mass bins, 95~\% bands). Top: the first three accepted attempts under four aggregations. Bottom: the answer of the stopping rule ($p_{\mathrm{stop}} = 0.90$), and its ECE against the perfect-calibration null band (right; * production). Only $s_k$ fitted on sessions is consistent with calibration at the stop, at the edge of its band.}
  \label{fig:M4_06_k3_calibration}\label{fig:app:prod-k3cal}
\end{figure}
\gentab{m4/m4_stop_calibration}
\gentab{m4/m4_stopping}
\FloatBarrier

% -------------------------------------------------------------------------------------------------
\subsection{Screen away and the rest-pose grip check}\label{app:prod:grip}

\textbf{Problem and detection.}
Every held-out recording was also scored as its screen-away twin: the phone turned 180 degrees about its long axis, so the device $x$ and $z$ axes change sign.
The gravity-frame channels and the yaw8 input do not change, but the strongest segment features are device-frame quantities along the screen normal (Appendix~\ref{app:prod:detector}).
Screen away, the production pipeline falls to \Rci{m4/stressflip_production_rec_acc}; it scores the return in \R{m4/stressflip_production_rec_ret} of the twins, and its answers are right \R{m4/stressflip_production_rec_acc_answered} of the time (Figure~\ref{fig:app:prod-grip}).
A confident wrong percentage is the worst outcome for the product.

\textbf{Options.}
(1) A check of the chosen segment's own start pose; (2) the grip-normalised feature set, invariant by construction; (3) a model-free check of the rest pose on the inference page: while the phone rests at arm's length before the first movement, the screen-normal component of the up vector must be at least \R{m4/grip_gate_upz}; otherwise the page asks to turn the screen towards the face and does not answer (scored as an abstention).

\textbf{Evidence.}
(1) The segment check stopped \R{m4/stressflip_production_rec_gate_seg} of the twins and left their accuracy at \R{m4/stressflip_production_rec_acc_gated_seg}; the answers it let through were right \R{m4/stressflip_production_rec_acc_answered_gated_seg} of the time.
When the detector takes the return for the approach, the chosen segment starts at the nose, so its start pose says nothing about the grip.
(2) The grip set scores \R{m4/stressflip_grip_uncal_rec_acc} screen away and \R{m4/grip_acc} as recorded, but grip minus robust on the recordings is \Rcib{m4/d_grip_vs_robust}~points (Holm $p = \R{m4/d_grip_vs_robust_p_holm}$, family \emph{trial}).
(3) The rest-pose check stops \R{m4/stressflip_production_rec_gate} of the twins.
The share answered wrongly falls from \R{m4/gripwrongflip_rec_without} to \R{m4/gripwrongflip_rec_with} (\Rb{m4/gripwrongflip_rec_d}~[\Rb{m4/gripwrongflip_rec_lo}, \Rb{m4/gripwrongflip_rec_hi}]~points, Holm $p < 0.001$, family \emph{grip}); the twins' half-credit accuracy moves to \Rci{m4/stressflip_production_rec_acc_gated}, because a stopped recording scores as an abstention.
On the screen-toward recordings it stops \R{m4/stress_production_rec_gate}: coverage goes from \R{m4/stress_production_rec_cov} to \R{m4/stress_production_rec_cov_gated}, the share answered wrongly from \R{m4/gripwrong_rec_without} to \R{m4/gripwrong_rec_with}, and the accuracy from \R{m4/stress_production_rec_acc} to \R{m4/stress_production_rec_acc_gated} (paired \Rb{m4/stresspair_grip_check_rec_d}~[\Rb{m4/stresspair_grip_check_rec_lo}, \Rb{m4/stresspair_grip_check_rec_hi}]~points, Holm $p = \R{m4/stresspair_grip_check_rec_p_holm}$, family \emph{stress}).

\textbf{Decision and rationale.}
Production keeps the robust features behind the rest-pose grip check, and the instructions ask for the screen facing the user, tilted slightly upwards.
A user asked to turn the phone then gets the screen-toward accuracy.
The first remedy was replaced; the grip set was rejected for its cost on real grips.
Caveats: the check reads gravity only, so a phone held near vertical or tilted forward is also asked to adjust; it was evaluated on synthetic twins, not on real screen-away recordings; only the recorded grips (upright, top left and top right, all screen toward) were tested.

\begin{figure}[htbp]
  \centering
  \genfig{m4/M4_12_screen_flip}
  \caption{Every held-out recording scored as recorded (screen toward) and as its screen-away twin, by the pipeline of its fold (half-credit accuracy, 95~\% CI). Hatched: with the rest-pose grip check, which stops \R{m4/stressflip_production_rec_gate} of the twins and \R{m4/stress_production_rec_gate} of the recordings. Supports the grip check over the grip-normalised features.}
  \label{fig:M4_12_screen_flip}\label{fig:app:prod-grip}
\end{figure}
\FloatBarrier

% -------------------------------------------------------------------------------------------------
\subsection{Platform and condition}\label{app:prod:breakdown}

Table~\ref{tab:m4_breakdown} stratifies the production pipeline by platform and by instructed condition; the conditions stratify only and never enter a model.
The platforms score alike: Android \Rci{m4/cal_nested_android_acc} on \R{m4/cal_nested_android_n} recordings, iOS \Rci{m4/cal_nested_ios_acc} on \R{m4/cal_nested_ios_n}.
Fast movements score higher than natural and slow ones, and the sideways grips higher than the upright grip; the slow, tilted-touch, twist and curved variants score lowest, with wide intervals.
The bent-start variant has fewer than 5 persons and no CI.
These strata are descriptive: no comparison between them was tested here (lane sec tests speed and grip; Appendix~\ref{app:sec}).

\gentab{m4/m4_breakdown}
\FloatBarrier

% -------------------------------------------------------------------------------------------------
\subsection{A failure mode: natural movements are rejected}\label{app:prod:natural}

\textbf{Problem.}
The first free-choice movement of each person is the most production-like trial in the data: unpractised and chosen freely.
The production detector rejects \R{m4/nat_rejected} of these \R{m4/nat_n} natural-use recordings (\R{m4/nat_rejection_rate}), against \R{m4/prim_rejection_rate} of the primary recordings.
From each person's own primary rejection rate, \R{m4/nat_expected_rejections} rejections are expected (Poisson-binomial $p = \R{m4/nat_p_poisson}$; Fisher's exact test, natural against primary, $p = \R{m4/nat_p_fisher}$; post hoc tests, unadjusted).

\textbf{What the pipeline achieves on natural use.}
It answers \R{m4/prod_nat_answered} of \R{m4/prod_nat_n} (\R{m4/prod_nat_coverage}; \R{m4/prod_nat_answered_tie} of them at exactly 0.5), with a right-hand recall of \R{m4/prod_nat_acc_answered} on the answered trials (\R{derived/m4_prod_nat_shown_acc} on the \R{derived/m4_prod_nat_n_shown} with a hand shown), \R{m4/prod_nat_acc_forced} in forced choice and \Rci{m4/prod_nat_acc} with half credit; the oracle window gives \R{m4/oracle_nat_acc}.
All free-choice trials are right-hand so far, so this is right-hand recall: a constant right answer would score 100~\%.

\textbf{Cause (post hoc).}
A candidate segment is the stored approach in \R{m4/nat_approach_found} of the natural recordings (\R{m4/prim_approach_found} primary), and that candidate's P(approach) is below the threshold in \R{m4/nat_approach_p_below} of them (\R{m4/prim_approach_p_below} primary).
The natural approaches move less along the screen normal: median \R{m4/natfeat_disp_z0_nat}~m against \R{m4/natfeat_disp_z0_prim}~m, with \R{m4/natfeat_disp_z0_below} of them below the primary 25th percentile; their median P(approach) is \R{m4/natfeat_p_approach_nat} against \R{m4/natfeat_p_approach_prim}.
The grip check fires on \R{m4/nat_grip_gate} of the natural and \R{m4/prim_grip_gate} of the primary recordings, so the grip does not explain the rejections.

\textbf{Consequence.}
A first, unpractised movement is rejected as ``not an approach'' more often than a practised one.
In production the page then asks for another touch; whether practised cue-free touches behave like the primary trials needs real cue-free streams.

% -------------------------------------------------------------------------------------------------
\subsection{Online operation, caveats and next steps}\label{app:prod:caveats}

\textbf{What runs on the phone.}
(1) Every standardised 60~Hz sample updates the causal activity.
(2) Before the first touch, the rest-pose grip check runs.
(3) Bouts are cut by hysteresis; a bout is a movement after \R{m4/seg_sustain_ms}~ms above \R{m4/seg_hi_dps}~deg/s.
(4) An approach ends at the first stillness at the face and is confirmed \R{m4/lookahead_s}~s later; its window is cut by the same code as the training windows.
(5) The segment classifier drops returns and other movements; every accepted approach adds its calibrated logit.
(6) The stopping rule of Appendix~\ref{app:prod:stopping} decides when to answer.
The same segmentation code runs online and offline; the replay of Appendix~\ref{app:prod:pipelines} checks it on re-standardised raw records.
The page must also check the instructed 1~s of stillness before it accepts a touch and say ``hold still'' otherwise.

\textbf{Caveats.}
\begin{itemize}[nosep]
  \item The production numbers use yaw8, lane m4's reference hand classifier, not the final model.
  \item The recordings come from a protocol with a GO cue and an enforced stillness; the stress test and the synthetic sessions approximate production but are spliced from those recordings, and the session-fitted scale is learned from such sessions too.
  \item Exploratory choices on the outer folds: the detector feature set (\R{m4/n_configs_detector} sets), the hand model, the attempt rule and the aggregation (\R{m4/n_configs_sessions} session configurations); the production result records \R{m4/n_configs_final} cue-free pipelines. The segmentation setting and the calibrator are nested.
  \item Mild in-fold optimism, not leakage of the held-out person: the recording-fitted scale and the calibrator are fitted on the same cross-fitted logits, and the session-fitted scale uses the outer calibrator on the same training persons.
  \item Untested: iOS sample stalls during an approach; a shorter hold at the nose (the stillness after a segment is a feature); real screen-away and landscape grips.
  \item Session-level intervals are wide (\R{m4/sess_n} sessions of \R{m4/n_persons} persons); natural use has one right-hand trial per person.
\end{itemize}

\textbf{Next steps.}
\begin{itemize}[nosep]
  \item Run the final model inside the production simulation (the hand model is a plug-in) and report its cue-free and multi-touch numbers.
  \item Reject the arm settling before an approach and returns, for example with a sequence model over the bouts (rest, approach, hold, return); the true-approach curve of Figure~\ref{fig:app:prod-attempts} is its ceiling.
  \item Add the return head to the stopping rule, with its own cross-fitted calibration.
  \item Augment the training streams with short stillness and screen-away grips; choose the P(approach) threshold inside the folds.
  \item Record a rehearsal session of real cue-free streams, with several touches per person, and re-test the natural-use rejections on it.
\end{itemize}
